\section{AMC 10 10A 2008}
        \begin{enumerate}
        	\item (\textit{AMC 10 10A 2008}) A bakery owner turns on his doughnut machine at $\text{8:30}\ {\small\text{AM}}$. At $\text{11:10}\ {\small\text{AM}}$ the machine has completed one third of the day's job. At what time will the doughnut machine complete the job?


 $\mathrm{(A)}\ \text{1:50}\ {\small\text{PM}}\qquad\mathrm{(B)}\ \text{3:00}\ {\small\text{PM}}\qquad\mathrm{(C)}\ \text{3:30}\ {\small\text{PM}}\qquad\mathrm{(D)}\ \text{4:30}\ {\small\text{PM}}\qquad\mathrm{(E)}\ \text{5:50}\ {\small\text{PM}}$


 
 \textbf{Answer: }D



\textbf{Solution}

The machine completes one-third of the job in $\text{11:10}-\text{8:30}=\text{2:40}$ hours. Thus, the entire job is completed in $3\cdot(\text{2:40})=\text{8:00}$ hours.

 Since the machine was started at $\text{8:30 AM}$, the job will be finished $8$ hours later, at $\text{4:30 PM}$. The answer is $\mathrm{(D)}$.

 Note: $\text{2:40}$ means $2$ hours and $40$ minutes. $3$ multiplied by this time interval is $8$ hours.

 


	\item (\textit{AMC 10 10A 2008}) A square is drawn inside a rectangle. The ratio of the width of the rectangle to a side of the square is $2:1$. The ratio of the rectangle's length to its width is $2:1$. What percent of the rectangle's area is in the square?


 $\mathrm{(A)}\ 12.5\qquad\mathrm{(B)}\ 25\qquad\mathrm{(C)}\ 50\qquad\mathrm{(D)}\ 75\qquad\mathrm{(E)}\ 87.5$


 
 \textbf{Answer: }A



\textbf{Solution 1}

Since they are asking for the "ratio" of two things, we can say that the side of the square is anything that we want. So if we say that it is $1$, then the width of the rectangle is $2$, and the length is $4$, thus making the total area of the rectangle $8$. The area of the square is just $1$. So the percent of the rectangle's area inside the square is $\frac{1}{8} \times 100 = 12.5 \longrightarrow \fbox{A}$

 


	\item (\textit{AMC 10 10A 2008}) For the positive integer $n$, let $\langle n\rangle$ denote the sum of all the positive divisors of $n$ with the exception of $n$ itself. For example, $\langle 4\rangle=1+2=3$ and $\langle 12 \rangle =1+2+3+4+6=16$. What is $\langle\langle\langle 6\rangle\rangle\rangle$?


 $\mathrm{(A)}\ 6\qquad\mathrm{(B)}\ 12\qquad\mathrm{(C)}\ 24\qquad\mathrm{(D)}\ 32\qquad\mathrm{(E)}\ 36$


 
 \textbf{Answer: }A



\textbf{Solution 1}

$\langle\langle\langle 6\rangle\rangle\rangle= \langle\langle 6\rangle\rangle = \langle 6\rangle=\ 6\quad\Longrightarrow\quad\mathrm{(A)}$

 



\textbf{Solution 2}

Since $6$ is a perfect number, any such operation where $n=6$ will yield $6$ as the answer. 

 Note: A perfect number is defined as a number that equals the sum of its positive divisors excluding itself.

 


	\item (\textit{AMC 10 10A 2008}) Suppose that $\frac{2}{3}$ of $10$ bananas are worth as much as $8$ oranges. How many oranges are worth as much as $\frac{1}{2}$ of $5$ bananas?


 $\mathrm{(A)}\ 2\qquad\mathrm{(B)}\ \frac{5}{2}\qquad\mathrm{(C)}\ 3\qquad\mathrm{(D)}\ \frac{7}{2}\qquad\mathrm{(E)}\ 4$


 
 \textbf{Answer: }C



\textbf{Solution}

If $\frac{2}{3}\cdot10\ \text{bananas}=8\ \text{oranges}$, then $\frac{1}{2}\cdot5\ \text{bananas}=\left(\frac{1}{2}\cdot 5\ \text{bananas}\right)\cdot\left(\frac{8\ \text{oranges}}{\frac{2}{3}\cdot10\ \text{bananas}}\right)=3\ \text{oranges}\Longrightarrow\mathrm{(C)}$. 

 


	\item (\textit{AMC 10 10A 2008}) Which of the following is equal to the product

 \[\frac{8}{4}\cdot\frac{12}{8}\cdot\frac{16}{12}\cdot\cdots\cdot\frac{4n+4}{4n}\cdot\cdots\cdot\frac{2008}{2004}?\]


 $\mathrm{(A)}\ 251\qquad\mathrm{(B)}\ 502\qquad\mathrm{(C)}\ 1004\qquad\mathrm{(D)}\ 2008\qquad\mathrm{(E)}\ 4016$


 
 \textbf{Answer: }B



\textbf{Solution 1}

$\frac {8}{4}\cdot\frac {12}{8}\cdot\frac {16}{12}\cdots\frac {4n + 4}{4n}\cdots\frac {2008}{2004} = \frac {1}{4}\cdot\left(\frac {8}{8}\cdot\frac {12}{12}\cdots\frac {4n}{4n}\cdots\frac {2004}{2004}\right)\cdot 2008 = \frac{2008}{4} =$ $502 \Rightarrow B$. 

 



\textbf{Solution 2}

Notice that everything cancels out except for $2008$ in the numerator and $4$ in the denominator.

 Thus, the product is $\frac{2008}{4}=502$, and the answer is $\textbf{(B)}$. 

 


	\item (\textit{AMC 10 10A 2008}) A triathlete competes in a triathlon in which the swimming, biking, and running segments are all of the same length. The triathlete swims at a rate of 3 kilometers per hour, bikes at a rate of 20 kilometers per hour, and runs at a rate of 10 kilometers per hour. Which of the following is closest to the triathlete's average speed, in kilometers per hour, for the entire race?


 $\mathrm{(A)}\ 3\qquad\mathrm{(B)}\ 4\qquad\mathrm{(C)}\ 5\qquad\mathrm{(D)}\ 6\qquad\mathrm{(E)}\ 7$


 
 \textbf{Answer: }D



\textbf{Solution}

Let $d$ be the length of one segment of the race.

 Average speed is total distance divided by total time. The total distance is $3d$, and the total time is $\frac{d}{3}+\frac{d}{20}+\frac{d}{10}=\frac{29d}{60}$.

 Thus, the average speed is $3d\div\left(\frac{29d}{60}\right)=\frac{180}{29}$. This is closest to $6$, so the answer is $\mathrm{(D)}$.

 



\textbf{Solution 2}

Since the three segments are all the same length, the triathlete's average speed is the \href{https://artofproblemsolving.com/wiki/index.php?title=Harmonic\_mean}{harmonic mean} of the three given rates. Therefore, the average speed is 
 \[\frac{3}{\frac{1}{3}+\frac{1}{20}+\frac{1}{10}}=\frac{3}{\frac{29}{60}}=\frac{180}{29}\approx6\Rightarrow\boxed{\mathrm{(D)}\ 6}\].

 


	\item (\textit{AMC 10 10A 2008}) The fraction

 \[\frac{\left(3^{2008}\right)^2-\left(3^{2006}\right)^2}{\left(3^{2007}\right)^2-\left(3^{2005}\right)^2}\]
simplifies to which of the following?


 $\mathrm{(A)}\ 1\qquad\mathrm{(B)}\ \frac{9}{4}\qquad\mathrm{(C)}\ 3\qquad\mathrm{(D)}\ \frac{9}{2}\qquad\mathrm{(E)}\ 9$


 
 \textbf{Answer: }E



\textbf{Solution}

Simplifying, we get 
 \[\frac{3^{4016}-3^{4012}}{3^{4014}-3^{4010}}.\] Factoring out $3^{4012}$ in the numerator and factoring out $3^{4010}$ in the denominator gives us 
 \[\frac{(3^4-1)(3^{4012})}{(3^4-1)(3^{4010})}.\] Canceling out $3^4-1$ gives us $\frac{3^{4012}}{3^{4010}}=\frac{3^2}{3^0}=9\ \mathrm{(E)}.$

 



\textbf{Solution 2}

Using Difference of Squares, $\frac{(3^{2008})^{2}-(3^{2006})^{2}}{(3^{2007})^{2}-(3^{2005}){^2}}$ becomes

 $\frac{(3^{2008}+3^{2006})(3^{2008}-3^{2006})}{(3^{2007}+3^{2005})(3^{2007}-3^{2005})}$

 
$= \frac{3^{2006}(9+1) \cdot 3^{2006}(9-1)}{3^{2005}(9+1) \cdot 3^{2005}(9-1)}$

 $= \boxed{\text{(E)}9}$

 


	\item (\textit{AMC 10 10A 2008}) Heather compares the price of a new computer at two different stores. Store $A$ offers $15\%$ off the sticker price followed by a $\textdollar90$ rebate, and store $B$ offers $25\%$ off the same sticker price with no rebate. Heather saves $\textdollar15$ by buying the computer at store $A$ instead of store $B$. What is the sticker price of the computer, in dollars?


 $\mathrm{(A)}\ 750\qquad\mathrm{(B)}\ 900\qquad\mathrm{(C)}\ 1000\qquad\mathrm{(D)}\ 1050\qquad\mathrm{(E)}\ 1500$


 
 \textbf{Answer: }A



\textbf{Solution 1}

Let the sticker price be $x$.

 The price of the computer is $0.85x-90$ at store $A$ and $0.75x$ at store $B$.

 Heather saves $\$15$ at store $A$, so $0.85x-90+15=0.75x$.

 Solving, we find $x=750$, and thus the answer is $\mathrm{(A)}$.

 



\textbf{Solution 2}

The $\$ 90$ at store $A$ is $\$ 15$ greater than the additional $25\%-15\% = 10\%$ off at store $B$.

 Thus the $10\%$ off is equal to $\$ 90$ $-$ $\$ 15$ $=$ $\$ 75$, and therefore the sticker price is $\$ 750$.

 


	\item (\textit{AMC 10 10A 2008}) Suppose that

 \[\frac{2x}{3}-\frac{x}{6}\]
is an integer. Which of the following statements must be true about $x$?


 $\mathrm{(A)}\ \text{It is negative.}\\\qquad\mathrm{(B)}\ \text{It is even, but not necessarily a multiple of 3.}\\\qquad\mathrm{(C)}\ \text{It is a multiple of 3, but not necessarily even.}\\\qquad\mathrm{(D)}\ \text{It is a multiple of 6, but not necessarily a multiple of 12.}\\\qquad\mathrm{(E)}\ \text{It is a multiple of 12.}$


 
 \textbf{Answer: }B



\textbf{Solution}


 \[\frac{2x}{3}-\frac{x}{6}\quad\Longrightarrow\quad\frac{4x}{6}-\frac{x}{6}\quad\Longrightarrow\quad\frac{3x}{6}\quad\Longrightarrow\quad\frac{x}{2}\]For $\frac{x}{2}$ to be an integer, $x$ must be even, but not necessarily divisible by $3$. Thus, the answer is $\mathrm{(B)}$.

 


	\item (\textit{AMC 10 10A 2008}) Each of the sides of a square $S_1$ with area $16$ is bisected, and a smaller square $S_2$ is constructed using the bisection points as vertices. The same process is carried out on $S_2$ to construct an even smaller square $S_3$. What is the area of $S_3$?


 $\mathrm{(A)}\ \frac{1}{2}\qquad\mathrm{(B)}\ 1\qquad\mathrm{(C)}\ 2\qquad\mathrm{(D)}\ 3\qquad\mathrm{(E)}\ 4$


 
 \textbf{Answer: }E



\textbf{Solution 1}

Since the area of the large square is $16$, the side length equals $4$. If all sides are bisected, the resulting square has side length $2\sqrt{2}$, thus making the area $8$. If we repeat this process again, we notice that the area is just half that of the previous square, so the area of $S_{3} = 4 \longrightarrow \fbox{E}$

 



\textbf{Solution 2}

Since the length ratio is $\frac{1}{\sqrt{2}}$, then the area ratio is $\frac{1}{2}$ (since the area ratio between two similar 2-dimensional objects is equal to the side ratio of those objects squared). This means that $S_2 = 8$ and $S_3 = \boxed{\textbf{(E) }4}$.

 


	\item (\textit{AMC 10 10A 2008}) While Steve and LeRoy are fishing 1 mile from shore, their boat springs a leak, and water comes in at a constant rate of 10 gallons per minute. The boat will sink if it takes in more than 30 gallons of water. Steve starts rowing toward the shore at a constant rate of 4 miles per hour while LeRoy bails water out of the boat. What is the slowest rate, in gallons per minute, at which LeRoy can bail if they are to reach the shore without sinking?


 $\mathrm{(A)}\ 2\qquad\mathrm{(B)}\ 4\qquad\mathrm{(C)}\ 6\qquad\mathrm{(D)}\ 8\qquad\mathrm{(E)}\ 10$


 
 \textbf{Answer: }D



\textbf{Solution 1}

It will take $\frac{1}{4}$ of an hour or $15$ minutes to get to shore.

 Since only $30$ gallons of water can enter the boat, only $\frac{30}{15}=2$ net gallons can enter the boat per minute.

 Since $10$ gallons of water enter the boat each minute, LeRoy must bail $10-2=8$ gallons per minute $\Rightarrow\mathrm{(D)}$.

 


 



\textbf{Solution 2}

We set up the following equation, where $x$ is the answer:

 $\frac{30}{10-x} = 15 \Rightarrow\mathrm{x=8}$, so the answer is $(D)$.

 



\textbf{Solution 3}

Let $r$ be the slowest rate LeRoy can bail out water. It will take them $15$ minutes to get to shore. As stated in solution 1, the number of gallons that enter the boat can be at most $2$ for the boat to stay afloat. Water enters at a rate of $10$ gallons per minute. 

 Therefore:

 $2 \geq 10-r$

 $r \geq 8$

 The minimum possible value for $r$ is $8\ \mathrm{(D)}$.

 


	\item (\textit{AMC 10 10A 2008}) In a collection of red, blue, and green marbles, there are $25\%$ more red marbles than blue marbles, and there are $60\%$ more green marbles than red marbles. Suppose that there are $r$ red marbles. What is the total number of marbles in the collection?
$\mathrm{(A)}\ 2.85r\qquad\mathrm{(B)}\ 3r\qquad\mathrm{(C)}\ 3.4r\qquad\mathrm{(D)}\ 3.85r\qquad\mathrm{(E)}\ 4.25r$


 
 \textbf{Answer: }C



\textbf{Solution}

The number of blue marbles is $\frac{4}{5}r$, the number of green marbles is $\frac{8}{5}r$, and the number of red marbles is $r$.

 Thus, the total number of marbles is $\frac{4}{5}r+\frac{8}{5}r+r=3.4r$, and the answer is $\mathrm{(C)}$.

 



\textbf{Solution 2}

Let the number of red marbles be 100. You have $b = 80,$ and $g = 160.$ 

 The answer is $\frac{100+80+160}{100} = \boxed{3.4}$

 


	\item (\textit{AMC 10 10A 2008}) Doug can paint a room in $5$ hours. Dave can paint the same room in $7$ hours. Doug and Dave paint the room together and take a one-hour break for lunch. Let $t$ be the total time, in hours, required for them to complete the job working together, including lunch. Which of the following equations is satisfied by $t$?


 $\mathrm{(A)}\ \left(\frac{1}{5}+\frac{1}{7}\right)\left(t+1\right)=1\qquad\mathrm{(B)}\ \left(\frac{1}{5}+\frac{1}{7}\right)t+1=1\qquad\mathrm{(C)}\ \left(\frac{1}{5}+\frac{1}{7}\right)t=1\\\mathrm{(D)}\ \left(\frac{1}{5}+\frac{1}{7}\right)\left(t-1\right)=1\qquad\mathrm{(E)}\ \left(5+7\right)t=1$


 
 \textbf{Answer: }D



\textbf{Solution 1}

Doug can paint $\frac{1}{5}$ of a room per hour, Dave can paint $\frac{1}{7}$ of a room per hour, and the time they spend working together is $t-1$.

 Since rate multiplied by time gives output, $\left(\frac{1}{5}+\frac{1}{7}\right)\left(t-1\right)=1 \Rightarrow \mathrm{(D)}$

 



\textbf{Solution 2}

If one person does a job in $a$ hours and another person does a job in $b$ hours, the time it takes to do the job together is $\frac{ab}{a+b}$ hours.

 Since Doug paints a room in 5 hours and Dave paints a room in 7 hours, they both paint in $\frac{5\times7}{5+7} = \frac{35}{12}$ hours. They also take 1 hour for lunch, so the total time $t = \frac{35}{12} + 1$ hours.

 Looking at the answer choices, $(D)$ is the only one satisfied by $t = \frac{35}{12} + 1$.

 


	\item (\textit{AMC 10 10A 2008}) Older television screens have an aspect ratio of $4: 3$. That is, the ratio of the width to the height is $4: 3$. The aspect ratio of many movies is not $4: 3$, so they are sometimes shown on a television screen by "letterboxing" - darkening strips of equal height at the top and bottom of the screen, as shown. Suppose a movie has an aspect ratio of $2: 1$ and is shown on an older television screen with a $27$-inch diagonal. What is the height, in inches, of each darkened strip?

                    \begin{figure}[H]
                        \centering
                        \includegraphics[width=0.5\textwidth]{Problem/AMC_10_10A_2008_14_Problem_Figure_1.png}
                    \end{figure}
                    
$\mathrm{(A)}\ 2\qquad\mathrm{(B)}\ 2.25\qquad\mathrm{(C)}\ 2.5\qquad\mathrm{(D)}\ 2.7\qquad\mathrm{(E)}\ 3$


 
 \textbf{Answer: }D



\textbf{Solution}

Let the width and height of the screen be $4x$ and $3x$ respectively, and let the width and height of the movie be $2y$ and $y$ respectively. 

 By the \href{https://artofproblemsolving.com/wiki/index.php?title=Pythagorean\_Theorem}{Pythagorean Theorem}, the diagonal is $\sqrt{(3x)^2+(4x)^2}=5x = 27$. So $x=\frac{27}{5}$. 

 Since the movie and the screen have the same width, $2y=4x\Rightarrow y=2x$. 

 Thus, the height of each strip is $\frac{3x-y}{2}=\frac{3x-2x}{2}=\frac{x}{2}=\frac{27}{10}=2.7\Longrightarrow\mathrm{(D)}$.

 


 


 



\textbf{Solution 2}

If the ratio is 4:3, and the diagonal is 27, two sides must have been multiplied by a factor of 27/5 to result in the diagonal being 27 (Since 3,4,5 forms a right triangle), multiplying by 27/5, we get width=108/5, height=81/5. From there, we can set up an equation, as it states that after the letterbox has been applied, the aspect ratio became 2:1. 

 We observe that the aspect ratio only takes away from the screen's height, not width. So, if we call X, the height of one darkening strip, we can write this equation: 
 \[\frac{\frac{108}{5}}{\frac{81}{5} - 2x} = 2\]. 

 Therefore, $x = \frac{27}{10} = 2.7$.

 
Thus, the answer is $\text{\textbf{(D)}}$.

 


	\item (\textit{AMC 10 10A 2008}) Yesterday Han drove 1 hour longer than Ian at an average speed 5 miles per hour faster than Ian. Jan drove 2 hours longer than Ian at an average speed 10 miles per hour faster than Ian. Han drove 70 miles more than Ian. How many more miles did Jan drive than Ian?


 $\mathrm{(A)}\ 120\qquad\mathrm{(B)}\ 130\qquad\mathrm{(C)}\ 140\qquad\mathrm{(D)}\ 150\qquad\mathrm{(E)}\ 160$


 
 \textbf{Answer: }D



\textbf{Solution 1}

Set the time Ian traveled as $I$, and set Han's speed as $H$. Therefore, Jan's speed is $H+5.$

 We get the following equation for how much Han is ahead of Ian:$H+5I = 70.$

 The expression for how much Jan is ahead of Ian is:$2(H+5)+10I.$

 This simplifies to: $2H+10+10I.$

 However, this is just $2(H+5I)+10.$

 Substitute, from the first equation, $H+5I$ as $70.$

 Therefore, the answer is $140 + 10$, which is $150$, or $\boxed{\mathrm{(D)}}$

 



\textbf{Solution 2}

We let Ian's speed and time equal $I_s$ and $I_t$, respectively. Similarly, let Han's and Jan's speed and time be $H_s$, $H_t$, $J_s$, $J_t$. The problem gives us 5 \href{https://artofproblemsolving.com/wiki/index.php?title=Equation}{equations}:

 
 \begin{align} H_s&=I_s+5 \\ H_t&=I_t+1 \\ J_s&=I_s+10 \\ J_t&=I_t+2 \\ H_s \cdot H_t & =I_s \cdot I_t+70 \end{align}

 Substituting equations $(1)$ and $(2)$ into $(5)$ gives:

 
 \[(I_s+5)(I_t+1)=I_s I_t+70 \Longrightarrow I_s I_t+I_s+5I_t+5=I_s I_t+70 \Longrightarrow I_s+5I_t=65 \quad (*)\]

 We are asked the difference between Jan's and Ian's distances, or

 
 \[J_s J_t-I_s I_t=x,\]

 Where $x$ is the difference between Jan's and Ian's distances and the answer to the problem. Substituting $(3)$ and $(4)$ equations into this equation gives:

 
 \[(I_s+10)(I_t+2)-I_s I_t=x \Longrightarrow I_s I_t+2I_s+10I_t+20-I_s I_t=x \Longrightarrow\]

 
 \[2I_s+10I_t+20=x \Longrightarrow 2(I_s+5I_t)+20=x\]

 Substituting $(*)$ into this equation gives:

 
 \[2(65)+20=x \Longrightarrow 130+20=x \Longrightarrow 150=x\]

 Therefore, the answer is $150$ miles or $\boxed{\mathrm{(D)}}$.

 



\textbf{Solution 3}

Let Ian drive $D$ miles, at a speed of $R$, for some time $T$ (in hours). Hence, we have $D=RT$. We can find a similar equation for Han, who drove $D + 70$ miles, at a rate of $R+5$, for $T+1$ hours, giving us $D + 70 = (R + 5)(T + 1)$. We can do the same for Jan, giving us $D + x = (R + 10)(T + 2)$, where $x$ is how much further Jan traveled than Ian. We now have three equations:
 \[D= RT\]
 \[D + 70 = (R+5)(T+1) = RT + R + 5T + 5\]
 \[D + x = (R + 10)(T + 2) = RT + 10 T + 2R + 20.\]Substituting $RT$ for $D$ in the second and third equations and cancelling gives us:
 \[70 = 5T + R + 5 \Longrightarrow 5T + R = 65\]
 \[x = 10T + 2R + 20  \Longrightarrow x = 2(5T + R ) + 20  \Longrightarrow x= 2(65) + 20 = 150.\]Since $x = 150$, our answer is $\boxed{\mathrm{(D)}}$.

 



\textbf{Solution 4}

Let Ian drive $d$ miles, $t$ hours, and at speed $s$.

 Ian's Equation: $d=s \cdot t.$

 Han drove 70 more miles, traveled 5 miles per hour faster and traveled 1 more hour than Ian.

 Han's Equation: $d+70=(s+5) \cdot (t+1).$

 Let Jan have driven $m$ miles. Jan also has driven 10 miles per hour faster and traveled 2 more hours than Ian.

 Jan's Equation: $m=(s+10) \cdot (t+2).$

 Let's group the equations together:

 $(1) \phantom{a} d=s \cdot t$

 $(2) \phantom{a} d+70=(s+5) \cdot (t+1)$

 $(3) \phantom{a} m=(s+10) \cdot (t+2)$

 Let's see what we want to find: We want to find $n$. The equation is $m=d+n$ where $n$ is the number of more miles traveled by Jan than Ian.

 Onto the calculating part.

 Expanding the second equation, we get

 $d+70=st+5t+s+5.$

 Note that $st=d$ by the first equation, so substituting we get

 $d+70=d+5t+s+5.$

 Simplifying gets us 

 $(4) \phantom{a} s+5t=65.$

 (Note for the above process: You could have substituted $d$ with $st$ but that would lead you to the same result since $d-d=st-st=0.$)

 Let's look at the third equation. Expanding, we get

 $m=st+10t+2s+20.$

 Since we want $n$ we want the equation $m=d+n$. We write the expanded third equation into this form since $st=d.$

 $(5) \phantom{a} m=d+(10t+2s+20)$.

 Let's take a closer look at the $n$ section of the equation:

 $(6) \phantom{a} n=10t+2s+20$

 This looks very similar to equation 4, if you multiply equation 4 by $2$ you get

 $(7) \phantom{a} 10t+2s=130$.

 Plugging equation 7 into equation 6 we have

 $n=(10t+2s)+20 \Rightarrow n=130+20$

 Calculating gets us

 $(8) \phantom{a} n=150.$

 Substituting equation 8 into equation 5 gets us

 $m=d+n \Rightarrow m=d+150.$

 The $n$ term is $150$ which is what we want to find so the answer is $\boxed{\mathrm{(D) 150}}.$

 



\textbf{Solution 5 (quick solution)}

Since Han drove for $1$ hour and drove $70$ miles more than Ian during that hour, we know that Ian's speed is $65$ miles per hour since Han drove $5$ mph faster than him. Now Jan went $10$ mph faster than Ian for $2$ hours, so we can tell that she drove $75 \cdot 2$ miles more than Ian, therefore the answer is $\boxed{\mathrm{(D) 150}}.$

 


	\item (\textit{AMC 10 10A 2008}) Points $A$ and $B$ lie on a circle centered at $O$, and $\angle AOB = 60^\circ$. A second circle is internally tangent to the first and tangent to both $\overline{OA}$ and $\overline{OB}$. What is the ratio of the area of the smaller circle to that of the larger circle?


 $\mathrm{(A)}\ \frac{1}{16}\qquad\mathrm{(B)}\ \frac{1}{9}\qquad\mathrm{(C)}\ \frac{1}{8}\qquad\mathrm{(D)}\ \frac{1}{6}\qquad\mathrm{(E)}\ \frac{1}{4}$


 
 \textbf{Answer: }B



\textbf{Solution 1}


                    \begin{figure}[H]
                        \centering
                        \includegraphics[width=0.5\textwidth]{Solution/AMC_10_10A_2008_16_Solution_Figure_1.png}
                    \end{figure}
                    Let $P$ be the center of the small circle with radius $r$, and let $Q$ be the point where the small circle is \href{https://artofproblemsolving.com/wiki/index.php?title=Tangent}{tangent} to $OA$. Also, let $C$ be the point where the small circle is tangent to the big circle with radius $R$.

 Then $PQO$ is a right triangle. Angle $POQ$ is $30$ degrees because line $OP$ bisects angle $AOB$ (this can be proved by dropping a perpendicular line from $P$ to line $OB$, letting their intersection be point $S$, and proving triangles $PQO$ and $PSO$ congruent), meaning that $PQO$ is a $30-60-90$ triangle. Therefore, $OP=2PQ$.

 Since $OP=OC-PC=OC-r=R-r$, we have $R-r=2PQ$, or $R-r=2r$, or $\frac{1}{3}=\frac{r}{R}$.

 Ratio of areas of circles is ratio of radii squared, so the answer is $\left(\frac{1}{3}\right)^2 = \frac{1}{9} \Rightarrow \boxed{B}$

 



\textbf{Solution 2}

Like in Solution 1, let $P$ be the center of the small circle with radius $r$, and let $Q$ be the point where the small circle is \href{https://artofproblemsolving.com/wiki/index.php?title=Tangent}{tangent} to $OA$.

 Let $N$ be the tangency point of the two circles. As shown in Solution 1, $POQ = 30$ degrees, so angle $NOA$ is also $30$ degrees. Let the line tangent to the two circles at $N$ intersect $OA$ and $OB$ at points $C$ and $D$, respectively. Since line $CD$ is tangent to circles $O$ and $P$, it must be perpendicular to $ON$, meaning that angle $ONC$ must be $90$ degrees. Because angle $NOA$ is $30$ degrees, angle $DCO$ is $180-30-90 = 60$ degrees. Angle $DOC$ is also $60$ degrees, so triangle $DOC$ is equilateral.

 Note that an equilateral triangle's incenter is also its centroid. This means the center of the inscribed circle is also the centroid. From properties of median lengths, the radius of the large circle is 3 times the radius of the small circle. 

 Then the ratio of areas will be $\frac{1}{3}$ squared, or $\frac{1}{9}\Rightarrow \boxed{\text{B}}$.

 


	\item (\textit{AMC 10 10A 2008}) An equilateral triangle has side length 6. What is the area of the region containing all points that are outside the triangle but not more than 3 units from a point on the triangle?


 $\mathrm{(A)}\ 36+24\sqrt{3}\qquad\mathrm{(B)}\ 54+9\pi\qquad\mathrm{(C)}\ 54+18\sqrt{3}+6\pi\qquad\mathrm{(D)}\ \left(2\sqrt{3}+3\right)^2\pi\\\mathrm{(E)}\ 9\left(\sqrt{3}+1\right)^2\pi$


 
 \textbf{Answer: }B



\textbf{Solution 1}


                    \begin{figure}[H]
                        \centering
                        \includegraphics[width=0.5\textwidth]{Solution/AMC_10_10A_2008_17_Solution_Figure_1.png}
                    \end{figure}
                    

 The region described contains three rectangles of dimensions $3 \times 6$, and three $120^{\circ}$ degree arcs of circles of \href{https://artofproblemsolving.com/wiki/index.php?title=Radius}{radius} $3$. Thus the answer is 
 \[3(3 \times 6) + 3 \left( \frac{120^{\circ}}{360^{\circ}} \times 3^2 \pi\right) = 54 + 9\pi \Longrightarrow \mathrm{(B)}.\]

 



\textbf{Solution 2}

After a quick sketch of the problem, one can deduce that there are $3$ rectangles in the figure, each with side length $6$ and width $3$. Therefore, the combined areas of the rectangles is $54$. The other three regions are circle-shaped areas, probably expressed in some form of $\pi$. Answer choices $\mathrm{(A)}$, $\mathrm{(D)}$, $\mathrm{(E)}$ are impossible because they either lack an integer or $\pi$ in the answer, and $\mathrm{(C)}$ is impossible since $18\sqrt{3}$ clearly does not belong to the rectangle or the circular areas. We can conclude that the only choice left is $\boxed{\mathrm{(B)}}$. 

 Solution by Airplane50

 


	\item (\textit{AMC 10 10A 2008}) A right triangle has perimeter 32 and area 20. What is the length of its hypotenuse?


 $\mathrm{(A)}\ \frac{57}{4}\qquad\mathrm{(B)}\ \frac{59}{4}\qquad\mathrm{(C)}\ \frac{61}{4}\qquad\mathrm{(D)}\ \frac{63}{4}\qquad\mathrm{(E)}\ \frac{65}{4}$


 
 \textbf{Answer: }B



\textbf{Solution 1}

Let the legs of the triangle have lengths $a,b$. Then, by the \href{https://artofproblemsolving.com/wiki/index.php?title=Pythagorean\_Theorem}{Pythagorean Theorem}, the length of the hypotenuse is $\sqrt{a^2+b^2}$, and the area of the triangle is $\frac 12 ab$. So we have the two equations

 $a+b+\sqrt{a^2+b^2} = 32 \\\\ \frac{1}{2}ab = 20$

 

 Re-arranging the first equation and squaring, 

 $\sqrt{a^2+b^2} = 32-(a+b)\\\\ a^2 + b^2 = 32^2 - 64(a+b) + (a+b)^2\\\\ a^2 + b^2 + 64(a+b) = a^2 + b^2 + 2ab + 32^2\\\\ a+b = \frac{2ab+32^2}{64}$

 

 From $(2)$ we have $2ab = 80$, so

 $a+b = \frac{80 + 32^2}{64} = \frac{69}{4}.$

 

 The length of the hypotenuse is $p - a - b = 32 - \frac{69}{4} = \frac{59}{4}\boxed{\ \mathrm{(B)}}$.

 



\textbf{Solution 2}

From the formula $A = rs$, where $A$ is the area of a triangle, $r$ is its \href{https://artofproblemsolving.com/wiki/index.php?title=Inradius}{inradius}, and $s$ is the \href{https://artofproblemsolving.com/wiki/index.php?title=Semiperimeter}{semiperimeter}, we can find that $r = \frac{20}{32/2} = \frac{5}{4}$. In a right triangle, $r = s - h$, where $h$ is the hypotenuse, so $h = 16 - \frac{5}{4} = \frac{59}{4}$ $\mathrm{(B)}$.

 \textbf{Note:} If the reader is unfamiliar with the inradius being equal to the semiperimeter minus the hypotenuse for a right triangle, it should not be too difficult to prove the relationship for themselves.~mobius247

 \textbf{Why $r=s-h$.} Draw a right triangle and inscribe a circle. Connect the center of the circle with tangency points; there are 3, one from each side of the triangle. Now, connect the center with each vertice and you will see 3 pairs of congruent triangles by HL (Hypotenuse-Leg Congruence). From these congruent triangles, mark equal lengths and you will find that $r=s-h$. 

 



\textbf{Solution 3}

From the problem, we know that 

 
 \begin{align*} a+b+c &= 32 \\ 2ab &= 80. \\ \end{align*}


 Subtracting $c$ from both sides of the first equation and squaring both sides, we get

 
 \begin{align*} (a+b)^2 &= (32 - c)^2\\ a^2 + b^2 + 2ab &= 32^2 + c^2 - 64c.\\ \end{align*}


 Now we substitute in $a^2 + b^2 = c^2$ as well as $2ab = 80$ into the equation to get

 
 \begin{align*} 80 &= 1024 - 64c\\ c &= \frac{944}{64}. \end{align*}


 Further simplification yields the result of $\frac{59}{4} \rightarrow \mathrm{(B)}$.

 



\textbf{Solution 4}

Let $a$ and $b$ be the legs of the triangle and $c$ the hypotenuse.

 Since the area is 20, we have $\frac{1}{2}ab = 20 => ab=40$.

 Since the perimeter is 32, we have $a + b + c = 32$.

 The Pythagorean Theorem gives $c^2 = a^2 + b^2$. 

 This gives us three equations with three variables:

 $ab = 40 \\ a + b + c = 32 \\ c^2 = a^2 + b^2$

 Rewrite equation 3 as $c^2 = (a+b)^2 - 2ab$.Substitute in equations 1 and 2 to get $c^2 = (32-c)^2 - 80$.

 $c^2 = (32-c)^2 - 80 \\\\ c^2 = 1024 - 64c + c^2 - 80 \\\\ 64c = 944 \\\\ c = \frac{944}{64} = \frac{236}{16} = \frac{59}{4}$. 

 The answer is choice $(B)$.

 



\textbf{Solution 5}

Let $a$, $b$, and $c$ be the sides of the triangle with $c$ as the hypotenuse.

 We know that $a + b + c =32$.

 According to the Pythagorean Theorem, we have $a^2 + b^2 = c^2$.

 We also know that $ab = 40$, since the area of the triangle is $20$.

 We substitute $2ab$ into $a^2 + b^2 = c^2$ to get $(a+b)^2 = c^2 + 80$.

 Moving the $c^2$ to the left, we again rewrite to get $(a+b+c)(a+b-c) = 80$.

 We substitute our value of $32$ for $a+b+c$ twice into our equation and subtract to get $a + b = \frac{69}{4}$.

 Finally, subtracting this from our original value of $32$, we get $\frac{59}{4}$, or $B$.

 



\textbf{Solution 6}

Let the legs be $a, b$. Then the hypotenuse is $\sqrt{a^2 + b^2}.$

 We know that $ab = 40$ and $a + b + \sqrt{a^2 + b^2} = 32$.

 The first equation gives 
 \[b = \frac{40}{a}\] and we can plug this into the second equation, yielding: 
 \[a + \frac{40}{a} + \sqrt{ a^2 + \frac{1600}{a^2}} = 32.\]

 Letting $X = a + \frac{40}{a}$, the equation becomes: 
 \[X + \sqrt{X^2 - 80} = 32.\]

 We can bring the $X$ to the right side and square which yields: 
 \[X^2 - 80 = (X - 32)^2 = X^2 - 64X + 1024.\]

 So, 
 \[-80 = -64X + 1024 \rightarrow X = \frac{69}{4}.\]

 Now, we know that $a + \frac{40}{a} = \frac{69}{4}.$

 Multiplying both sides by $4a$ gives: 
 \[4a^2 - 69a + 160.\]

 It can be observed that the roots of this equation are $a$ and $b$. We want the hypotenuse which is $\sqrt{a^2 + b^2} = \sqrt{ (a+b)^2 - 2ab}.$

 We can now apply Vieta's Formula which gives: 
 \[c = \sqrt{ \left( \frac{69}{4} \right)^2 - 80} = \boxed{ \frac{59}{4}}.\]

 



\textbf{Solution 7}

Let the sides be $a, b, c$ where $a$ and $b$ are the legs and $c$ is the hypotenuse.

 Since the perimeter is 32, we have

 $(1) \phantom{a} a+b+c=32$.

 Since the area is 20 and the legs are $a$ and $b$, we have that

 $(2) \phantom{a} \frac{a \cdot b}{2}=20$.

 By the Pythagorean Theorem, we have that

 $(3) \phantom{a} a^2+b^2=c^2$.

 Since we want $c$, we will equations $1, 2, 3$ be in the form of $c.$

 Equation 1 can be turned into

 $(4) \phantom{a} a+b=32-c$.

 Equation 2 can be simplified into

 $(5) \phantom{a} ab=40.$

 Equation 3 is already simplified.

 Onto the calculating process.

 Squaring the 1st equation we have

 $(a+b+c)^2=32^2.$

 Expanding and grouping, we have

 $(a^2+b^2+c^2)+2(ab+ac+bc)=32^2.$

 By equation 3 and substituting we get

 $2(c^2+ab+ac+bc)=32^2.$

 By equation 5 and substituting we get

 $2(c^2+40+ac+bc)=32^2.$

 Note that we can factor $c$ out in the inner expression, and we get

 $2(c^2+40+c(a+b))=32^2.$

 By equation 4 and substituting, we have

 $2(c^2+40+c(32-c))=32^2.$

 Expanding, we have

 $2(c^2+40+32c-c^2)=32^2.$

 Simplifying, we have

 $2(40+32c)=32^2.$

 Expanding again, we get

 $80+64c=32^2.$

 Dividing both sides by $16$ gets us

 $4c+5=64$

 Calculating gets us

 $c=\boxed{\mathrm{(B) \frac{59}{4}}}$.

 



\textbf{Solution 8}

This solution is very similar to Solution 1, except instead of subtracting $a+b$ from both sides in the first part, we subtract $\sqrt{a^2+b^2}$ from both sides, which gets us:

 
 \begin{align*} a+b&=32-\sqrt{a^2+b^2}\\ (a+b)^2&=(32-\sqrt{a^2+b^2})^2\\ a^2+2ab+b^2&=32^2-2\cdot32\cdot\sqrt{a^2+b^2}+a^2+b^2\\ 2ab&=32^2-2\cdot32\cdot\sqrt{a^2+b^2}.\\ \end{align*}


 We know that 
 \[\frac{1}{2}ab=20\]so 
 \[ab=40.\]

 We can then substitute $40$ for $ab$ to get us:

 
 \begin{align*} 80&=32^2-2\cdot32\cdot\sqrt{a^2+b^2}\\ 2\cdot32\cdot\sqrt{a^2+b^2}&=32^2-80\\ 64\cdot\sqrt{a^2+b^2}&=32^2-80\\ \sqrt{a^2+b^2}&=\frac{32^2-80}{64}\\ \end{align*}


 We know that $\sqrt{a^2+b^2}$ is the hypotenuse, so we only have to solve the right-hand side now.

 
 \begin{align*} \sqrt{a^2+b^2}&=\frac{32^2-80}{64}\\ &=\frac{32^2}{64}-\frac{80}{64}\\ &=\frac{32\cdot32}{32\cdot2}-\frac{32\cdot\frac{5}{2}}{32\cdot2}\\ &=\frac{32}{2}-\frac{\frac{5}{2}}{2}\\ &=16-\frac{5}{4}\\ &=\frac{64}{4}-\frac{5}{4}\\ &=\boxed{\mathrm{(B)} \frac{59}{4}}\\ \end{align*}


 


	\item (\textit{AMC 10 10A 2008}) Rectangle $PQRS$ lies in a plane with $PQ=RS=2$ and $QR=SP=6$. The rectangle is rotated $90^\circ$ clockwise about $R$, then rotated $90^\circ$ clockwise about the point $S$ moved to after the first rotation. What is the length of the path traveled by point $P$?


 $\mathrm{(A)}\ \left(2\sqrt{3}+\sqrt{5}\right)\pi\qquad\mathrm{(B)}\ 6\pi\qquad\mathrm{(C)}\ \left(3+\sqrt{10}\right)\pi\qquad\mathrm{(D)}\ \left(\sqrt{3}+2\sqrt{5}\right)\pi\\\mathrm{(E)}\ 2\sqrt{10}\pi$


 
 \textbf{Answer: }C



\textbf{Solution}


                    \begin{figure}[H]
                        \centering
                        \includegraphics[width=0.5\textwidth]{Solution/AMC_10_10A_2008_19_Solution_Figure_1.png}
                    \end{figure}
                    

 We let $P'Q'R'S'$ be the rectangle after the first rotation, and $P''Q''R''S''$ be the rectangle after the second rotation. Point $P$ pivots about $R$ in an \href{https://artofproblemsolving.com/wiki/index.php?title=Arc}{arc} of a circle of radius $\sqrt{2^2+6^2} = 2\sqrt{10}$, and since $\angle PRS,\, \angle P'RQ'$ are complementary, it follows that the arc has a degree measure of $90^{\circ}$ and length $\frac14$ of the \href{https://artofproblemsolving.com/wiki/index.php?title=Circumference}{circumference}. Thus, $P$ travels $\frac 14 \left(4\sqrt{10}\right)\pi = \sqrt{10}\pi$ in the first rotation.

 Similarly, in the second rotation, $P$ travels in a $90^{\circ}$ arc about $S'$, with the radius being $6$. It travels $\frac 14(12)\pi = 3\pi$. Therefore, the total distance it travels is $\left(3+\sqrt{10}\right)\pi\ \mathrm{(C)}$.

 


	\item (\textit{AMC 10 10A 2008}) Trapezoid $ABCD$ has bases $\overline{AB}$ and $\overline{CD}$ and diagonals intersecting at $K$. Suppose that $AB = 9$, $DC = 12$, and the area of $\triangle AKD$ is $24$. What is the area of trapezoid $ABCD$?


 $\mathrm{(A)}\ 92\qquad\mathrm{(B)}\ 94\qquad\mathrm{(C)}\ 96\qquad\mathrm{(D)}\ 98 \qquad\mathrm{(E)}\ 100$


 
 \textbf{Answer: }D



\textbf{Solution 1}


                    \begin{figure}[H]
                        \centering
                        \includegraphics[width=0.5\textwidth]{Solution/AMC_10_10A_2008_20_Solution_Figure_1.png}
                    \end{figure}
                    

 Since $\overline{AB} \parallel \overline{DC}$ it follows that $\triangle ABK \sim \triangle CDK$. Thus $\frac{KA}{KC} = \frac{KB}{KD} = \frac{AB}{DC} = \frac{3}{4}$. 

 We now introduce the concept of \href{https://artofproblemsolving.com/wiki/index.php?title=Area\_ratios}{area ratios}: given two triangles that share the same height, the ratio of the areas is equal to the ratio of their bases. Since $\triangle AKB, \triangle AKD$ share a common \href{https://artofproblemsolving.com/wiki/index.php?title=Altitude}{altitude} to $\overline{BD}$, it follows that (we let $[\triangle \ldots]$ denote the area of the triangle) $\frac{[\triangle AKB]}{[\triangle AKD]} = \frac{KB}{KD} = \frac{3}{4}$, so $[\triangle AKB] = \frac{3}{4}(24) = 18$. Similarly, we find $[\triangle DKC] = \frac{4}{3}(24) = 32$ and $[\triangle BKC] = 24$.

 Therefore, the area of $ABCD = [AKD] + [AKB] + [BKC] + [CKD] = 24 + 18 + 24 + 32 = 98\ \mathrm{(D)}$.

 



\textbf{Solution 2}

We may consider that trapezoid to be right, as there is nothing specifying its angles.Consider D and A right. Let the length of DA be h. Now we let A be (0,0) and we compute the x-coordinate of K from lines AC and DB. $y=\frac{h}{9}x$ for line DB, $y=-\frac{h}{12}x+h$ for line AC. Solving for K, $\frac{h}{9}x=-\frac{h}{12}x+h$ simplifying, $(\frac{1}{9}+\frac{1}{12})x=1$, $x=\frac{72}{14}=\frac{36}{7}$. Using the fact that $\frac{1}{2}*h*x=24$, we solve for h. $\frac{1}{2}*\frac{36}{7}*h=24, h=\frac{7}{18}*24=\frac{7*4}{3}$. Applying trapezoid area formula: $\frac{7*4}{3}*\frac{9+12}{2}=7*7*2=98$. Thus, the area is 98 and the answer is $\mathrm{(D)}$

 


	\item (\textit{AMC 10 10A 2008}) A cube with side length $1$ is sliced by a plane that passes through two diagonally opposite vertices $A$ and $C$ and the midpoints $B$ and $D$ of two opposite edges not containing $A$ or $C$, as shown. What is the area of quadrilateral $ABCD$?


 
                    \begin{figure}[H]
                        \centering
                        \includegraphics[width=0.5\textwidth]{Problem/AMC_10_10A_2008_21_Problem_Figure_1.png}
                    \end{figure}
                    


 $\mathrm{(A)}\ \frac{\sqrt{6}}{2}\qquad\mathrm{(B)}\ \frac{5}{4}\qquad\mathrm{(C)}\ \sqrt{2}\qquad\mathrm{(D)}\ \frac{5}{8}\qquad\mathrm{(E)}\ \frac{3}{4}$


 
 \textbf{Answer: }A



\textbf{Solution}


                    \begin{figure}[H]
                        \centering
                        \includegraphics[width=0.5\textwidth]{Solution/AMC_10_10A_2008_21_Solution_Figure_1.png}
                    \end{figure}
                    

 Since $AB = AD = CB = CD = \sqrt{\left(\frac{1}{2}\right)^2+1^2}$, it follows that $ABCD$ is a \href{https://artofproblemsolving.com/wiki/index.php?title=Rhombus}{rhombus}. The area of the rhombus can be computed by the formula $A = \frac 12 d_1d_2$, where $d_1,\,d_2$ are the diagonals of the rhombus (or of a \href{https://artofproblemsolving.com/wiki/index.php?title=Kite}{kite} in general). $BD$ has the same length as a face diagonal, or $\sqrt{1^2 + 1^2} = \sqrt{2}$. $AC$ is a space diagonal, with length $\sqrt{1^2+1^2+1^2} = \sqrt{3}$. Thus $A = \frac 12 \times \sqrt{2} \times \sqrt{3} = \frac{\sqrt{6}}{2}\ \mathrm{(A)}$.

 



\textbf{Solution 2}

We split $ABCD$ into $2$ congruent isosceles triangles. The height to the base $BD$ is clearly $\frac{\sqrt{3}}{2}$ because it is half of the space diagonal, and we can find the length of $BD$ to be $\sqrt{2}$ by the distance formula. 

 The area of each triangle is $\frac{\sqrt{3}}{2} \cdot \frac{\sqrt{2}}{2} = \frac{\sqrt{6}}{4}$. We multiply this by $2$ to find the area of $ABCD$ to be $\mathbf{\frac{\sqrt{6}}{2}}$.

 So, the answer is $\boxed{\text{A}}$.

 


	\item (\textit{AMC 10 10A 2008}) Jacob uses the following procedure to write down a sequence of numbers. First he chooses the first term to be 6. To generate each succeeding term, he flips a fair coin. If it comes up heads, he doubles the previous term and subtracts 1. If it comes up tails, he takes half of the previous term and subtracts 1. What is the probability that the fourth term in Jacob's sequence is an integer?


 $\mathrm{(A)}\ \frac{1}{6}\qquad\mathrm{(B)}\ \frac{1}{3}\qquad\mathrm{(C)}\ \frac{1}{2}\qquad\mathrm{(D)}\ \frac{5}{8}\qquad\mathrm{(E)}\ \frac{3}{4}$


 
 \textbf{Answer: }D



\textbf{Solution 1}

We construct a tree showing all possible outcomes that Jacob may get after $3$ flips; we can do this because there are only $8$ possibilities:
 \[6\quad\begin{cases}  \ \text{H}: 11 &\quad  \begin{cases}  \ \text{H}: 21 &\quad \begin{cases} \ \text{H}: \boxed{41}\\ \ \text{T}: 9.5 \end{cases}\\  \ \text{T}: 4.5 &\quad \begin{cases} \ \text{H}: \boxed{8}\\ \ \text{T}: 1.25 \end{cases}  \end{cases}\\  \ \text{T}: 2 &\quad  \begin{cases}  \ \text{H}: 3 &\qquad\! \begin{cases} \ \text{H}: \boxed{5}\\ \ \text{T}: 0.5 \end{cases}\\  \ \text{T}: 0 &\qquad\! \begin{cases} \ \text{H}: \boxed{-1}\\ \ \text{T}: \boxed{-1} \end{cases} \end{cases} \end{cases}\]There is a $\frac{5}{8}$ chance that the fourth term in Jacob's sequence is an integer, so the answer is $\mathrm{(D)}$.

 



\textbf{Solution 2}

We can see that as long as the last flip is heads, it will be an integer, so there is a $\dfrac{1}{2}$ chance of this happening.

 Doing a little casework, we see that the only possibility when the last flip is tails is when the third flip brings it to 0. There is a $\dfrac{1}{8}$ chance of this happening.

 Therefore, there is a $\dfrac{1}{2}+\dfrac{1}{8}=\boxed{(D)\dfrac{5}{8}}$ possibility of ending with an integer.

 


	\item (\textit{AMC 10 10A 2008}) Two subsets of the set $S=\lbrace a,b,c,d,e\rbrace$ are to be chosen so that their union is $S$ and their intersection contains exactly two elements. In how many ways can this be done, assuming that the order in which the subsets are chosen does not matter?


 $\mathrm{(A)}\ 20\qquad\mathrm{(B)}\ 40\qquad\mathrm{(C)}\ 60\qquad\mathrm{(D)}\ 160\qquad\mathrm{(E)}\ 320$


 
 \textbf{Answer: }B



\textbf{Solution 1}

First, choose the two letters to be repeated in each set.  $\dbinom{5}{2}=10$.  Now we have three remaining elements that we wish to place into two separate subsets. There are $2^3 = 8$ ways to do so because each of the three remaining letters can be placed either into the first or second subset. Both of those subsets contain the two chosen elements, so their intersection is the two chosen elements). Unfortunately, we have over-counted  (Take for example $S_{1} = \{a,b,c,d \}$ and $S_{2} = \{a,b,e \}$). Notice how $S_{1}$ and $S_{2}$ are interchangeable. Division by two will fix this problem. Thus we have:

 $\dfrac{10 \times 8}{2} = 40 \implies \boxed{\text{B}}$

 



\textbf{Solution 2}

Another way of looking at this problem is to break it down into cases.

 First, our two subsets can have 2 and 5 elements. The 5-element subset (aka the set) will contain the 2-element subset. There are $\dbinom{5}{2}=10$ ways to choose the 2-element subset. Thus, there are $10$ ways to create these sets.

 
Second, the subsets can have 3 and 4 elements. $3+4=7$ non-distinct elements. $7-5=2$ elements in the intersection. There are $\dbinom{5}{3}=10$ ways to choose the 3-element subset. For the 4-element subset, two of the elements must be the remaining elements (not in the 3-element subset). The other two have to be a subset of the 3-element subset. There are $\dbinom{3}{2}=3$ ways to choose these two elements, which means there are 3 ways to choose the 4-element subset. Therefore, there are $10\cdot3=30$ ways to choose these sets.

 
This leads us to the answer:

 $10+30=40 \implies \boxed{\textbf{(B) 40}}$

 



\textbf{Solution 3}

We label the subsets subset 1 and subset 2. Suppose the first subset has $k$ elements where $k<5.$ The second element has $5-k$ elements which the first subset does not contain (in order for the union to be the whole set). Additionally, the second set has 2 elements in common with the first subset. Therefore the number of ways to choose these sets is $\binom{5}{k} \cdot \binom{k}{2}.$ Computing for $k<5$ we have $10+30+30+10=80.$ Divide by 2 as order does not matter to get $40.$

 



\textbf{Solution 4 (Stars And Bars) [Isn't actually a valid solution]}

There are $\dbinom{5}{2}=10$ ways to choose the 2 shared elements. We now must place the 3 remaining elements into the subsets. Using stars and bars, we can notate this as: $I I I X \rightarrow \dbinom{4}{3}=4$. Thus, $4*10=\boxed{\textbf{(B) 40}}$

 Note from harryhero:

 This actually doesn't work for higher cases. In the 6-case, for example, we obtain an incorrect answer of $50$ using this method when it should be $80.$ This occurs due to the lack of order for the three elements - for example, having a, b on one side and d on another is different from b, c on one side and e on another, but we've counted both cases equally in the stars-and-bars method.

 Note from MaxIsW: This solution doesn't work for most cases and got lucky with the right numbers in this case.

 


	\item (\textit{AMC 10 10A 2008}) Let $k={2008}^{2}+{2}^{2008}$. What is the units digit of $k^2+2^k$?


 $\mathrm{(A)}\ 0\qquad\mathrm{(B)}\ 2\qquad\mathrm{(C)}\ 4\qquad\mathrm{(D)}\ 6\qquad\mathrm{(E)}\ 8$


 
 \textbf{Answer: }D



\textbf{Solution}

$k \equiv 2008^2 + 2^{2008} \equiv 8^2 + 2^4 \equiv 4+6 \equiv 0 \pmod{10}$. 

 So, $k^2 \equiv 0 \pmod{10}$. Since $k = 2008^2+2^{2008}$   is a multiple of four and the units digit of powers of two repeat in cycles of four, $2^k \equiv 2^4 \equiv 6 \pmod{10}$. 

 Therefore, $k^2+2^k \equiv 0+6 \equiv 6 \pmod{10}$. So the units digit is $6 \Rightarrow \boxed{D}$.

 \textbf{Note} Another way to get $k \equiv 0 \pmod{10}$ is to find the cycles of the last digit.

 For $2008^2$, we need only be concerned with the last digit $8$ since the other digits do not affect the last digit. Since $8^{2} = 64$, the last digit of $2008^2$ is $4$.

 For $2^{2008}$, note that the last digit cycles through the pattern ${2, 4, 8, 6}$. (You can try to see this by calculating the first powers of $2$.)

 Since $2008$ is a multiple of $4$, the last digit of $2^{2008}$ is evidently $6.$

 Continue as follows (more detail in Solution 2)

 



\textbf{Solution 2}

I am going to share another approach to this problem.

 A units digit $k$ for an integer $n$ implies $n \equiv k \pmod{10}$

 Let us take this step by step. First, we consider $k^2.$

 Note that $k^2 = \left(2008^2 + 2^{2008}\right)^2 = 2008^4 + 2 \cdot 2008^2 \cdot 2^{5016}.$ Now we calculate $k^2 \pmod{10}$

 Before continuing, though, we must take note of the following:

 
 \begin{align*} 2^1 &\equiv 2 \pmod {10} \\ 2^2 &\equiv 4 \pmod {10} \\ 2^3 &\equiv 8 \pmod {10} \\ 2^4 &\equiv 6 \pmod {10} \end{align*}


 Now, we continue with the calculation.

 
 \begin{align*} 2008^4 + 2 \cdot 2008^2 \cdot 2^{5016} &\equiv 8^4 + 2 \cdot 8^2 \cdot 2^{2008} + 2^{5016} \pmod{10} \\ &\equiv 8^4 + 2^1 \cdot 2^6 \cdot 2^{2008} + 2^{5016} \pmod{10} \\ &\equiv 8^4 + 2^{1+6+2008} + 2^{5016} \pmod{10} \\ &\equiv 6 + 2^{2015} + 6 \pmod{10} \\ &\equiv 6 + 2^{3} + 6 \pmod{10} \\ &\equiv 6 + 8 + 6 \pmod{10} \\ &\equiv 20 \pmod{10} \\ &\equiv 0 \pmod{10} \\ \end{align*}


 We do the same with $2^k.$ However, we just need to find $k \pmod 4$ in order to do this calculation since we have the table of $2^k \pmod 10.$

 
 \begin{align*} 2008^2 + 2^{2008} &\equiv 8^2 \pmod{4} \\ &\equiv 64 \pmod 4\\ &\equiv 0 \pmod 4 \end{align*}


 This implies that 
 \begin{align*} 2^k &\equiv 2^{4} \pmod{10} \\ &\equiv 6 \pmod{10} \end{align*}


 Thus, 
 \begin{align*} k^2 + 2^k &\equiv 6+0 \pmod{10}\\ &\equiv \boxed{\textbf{(D) }6} \pmod{10} \end{align*}


 



\textbf{Solution 3 (Dangerously Quick)}

By unit digit arithmetic, the unit digit of $k^2+2^k$ need to be either $4$ or $6.$ Hence, we can guess one of them for a $50$\% of getting it right.This should only take 20 seconds or less. ~peelybonehead

 Note: This solution is not recommended - this is effectively a truncated version of Solution 1 - do this only if you have very little time left.

 


	\item (\textit{AMC 10 10A 2008}) A round table has radius $4$. Six rectangular place mats are placed on the table. Each place mat has width $1$ and length $x$ as shown. They are positioned so that each mat has two corners on the edge of the table, these two corners being end points of the same side of length $x$. Furthermore, the mats are positioned so that the inner corners each touch an inner corner of an adjacent mat. What is $x$?


 
                    \begin{figure}[H]
                        \centering
                        \includegraphics[width=0.5\textwidth]{Problem/AMC_10_10A_2008_25_Problem_Figure_1.png}
                    \end{figure}
                    


 $\mathrm{(A)}\ 2\sqrt{5}-\sqrt{3}\qquad\mathrm{(B)}\ 3\qquad\mathrm{(C)}\ \frac{3\sqrt{7}-\sqrt{3}}{2}\qquad\mathrm{(D)}\ 2\sqrt{3}\qquad\mathrm{(E)}\ \frac{5+2\sqrt{3}}{2}$


 
 \textbf{Answer: }C



\textbf{Solution 1 (Trigonometry)}

Let one of the mats be $ABCD$, and the center be $O$ as shown: 

 
                    \begin{figure}[H]
                        \centering
                        \includegraphics[width=0.5\textwidth]{Solution/AMC_10_10A_2008_25_Solution_Figure_1.png}
                    \end{figure}
                    

 Since there are $6$ mats, $\Delta BOC$ is \href{https://artofproblemsolving.com/wiki/index.php?title=Equilateral}{equilateral} (the hexagon with side length $x$ is regular). So, $BC=CO=x$. Also, $\angle OCD = \angle OCB + \angle BCD = 60^\circ+90^\circ=150^\circ$. 

 By the \href{https://artofproblemsolving.com/wiki/index.php?title=Law\_of\_Cosines}{Law of Cosines}: $4^2=1^2+x^2-2\cdot1\cdot x \cdot \cos(150^\circ) \Rightarrow x^2 + x\sqrt{3} - 15 = 0 \Rightarrow x = \frac{-\sqrt{3}\pm 3\sqrt{7}}{2}$. 

 Since $x$ must be positive, $x = \frac{3\sqrt{7}-\sqrt{3}}{2} \Rightarrow (C)$.

 



\textbf{Solution 2 (without trigonometry)}

Draw $OD$ and $OC$ as in the diagram. Draw the altitude from $O$ to $DC$ and call the intersection $E$.

 

                    \begin{figure}[H]
                        \centering
                        \includegraphics[width=0.5\textwidth]{Solution/AMC_10_10A_2008_25_Solution_Figure_2.png}
                    \end{figure}
                    

 As proved in the first solution, $\angle OCD = 150^\circ$. That makes $\Delta OCE$ a $30-60-90$ triangle, so $OE = \frac{x}{2}$ and $CE= \frac{x\sqrt 3}{2}$

 Since $\Delta ODE$ is a right triangle, $\left({\frac{x}{2}}\right)^2 + \left({\frac{x\sqrt 3}{2} +1}\right)^2 = 4^2 \Rightarrow  x^2+x\sqrt3-15 = 0$

 Solving for $x$ gives $x =\frac{3\sqrt{7}-\sqrt{3}}{2}\Rightarrow (C)$

 



\textbf{Solution 3 (simply Pythagorean Theorem)}


                    \begin{figure}[H]
                        \centering
                        \includegraphics[width=0.5\textwidth]{Solution/AMC_10_10A_2008_25_Solution_Figure_3.png}
                    \end{figure}
                    

 By symmetry, $E$ is the midpoint of $DF$ and $OE$ is an extension of $OC$. Thus $\angle OED = 90^\circ$. Since $OD = 4$ and $DE = \frac{1}{2}$, $OE = \sqrt{16-\frac{1}{4}} = \frac{\sqrt{63}}{2} = \frac{3\sqrt{7}}{2}$. Since $\triangle CED$ is $30-60-90$, $CE = \frac{\sqrt{3}}{2}$ (this can also be deduced from Pythagoras on $\triangle CED$).

 Thus $OC = \frac{3\sqrt{7}-\sqrt{3}}{2}$. As previous solutions noted, $\triangle BOC$ is equilateral, and thus the desired length is $x = OC \implies (C)$.

 



\textbf{Solution 4}


                    \begin{figure}[H]
                        \centering
                        \includegraphics[width=0.5\textwidth]{Solution/AMC_10_10A_2008_25_Solution_Figure_4.png}
                    \end{figure}
                    

 Looking at the diagram above, we know that $BE$ is a diameter of circle $O$ due to symmetry. Due to \href{https://artofproblemsolving.com/wiki/index.php?title=Thales\%27\_theorem}{Thales' theorem}, triangle $ABE$ is a right triangle with $A = 90 ^\circ$. $AE$ lies on $AD$ and $GE$ because $BAD$ is also a right angle. To find the length of $DG$, notice that if we draw a line from $F$ to $M$, the midpoint of line $DG$, it creates two $30$ - $60$ - $90$ triangles. Therefore, $MD = MG = \frac{\sqrt{3}x}{2} \Rightarrow DG = \sqrt{3}x$. $AE = 2 + \sqrt{3}x$

 Applying the Pythagorean theorem to triangle $ABE$, we get 
 \[(2+\sqrt{3}x)^2 + x^2 = 8^2 \Rightarrow 4 + 3x^2 + 4\sqrt{3}x + x^2 = 64 \Rightarrow x^2 + \sqrt{3}x - 15 = 0\] Using the quadratic formula to solve, we get 
 \[x = \frac{-\sqrt{3} \pm \sqrt{3 -4(1)(-15)}}{2} = \frac{\pm 3\sqrt{7} - \sqrt{3}}{2}\] $x$ must be positive, therefore 
 \[x = \frac{3\sqrt{7} - \sqrt{3}}{2} \Rightarrow (C)\]

 



\textbf{Solution 5 (coordinate bashing)}


                    \begin{figure}[H]
                        \centering
                        \includegraphics[width=0.5\textwidth]{Solution/AMC_10_10A_2008_25_Solution_Figure_5.png}
                    \end{figure}
                    

 We will let $O(0,0)$ be the origin. This way the coordinates of $C$ will be $(0,y)$. By $30-60-90$, the coordinates of $D$ will be $\left(-\frac{1}{2}, y + \frac{\sqrt{3}}{2}\right)$. The distance any point with coordinates $(x, y)$ is from the origin is $\sqrt{x^2 + y^2}$. Therefore, the distance $D$ is from the origin is $4$ and $\frac{1}{4} + x^2 + x\sqrt{3} + \frac{3}{4} = x^2 + x\sqrt{3} + 1 = 16$. We get the quadratic equation mentioned in solution 2. Using the quadratic formula, we get that $x = \frac{3\sqrt{7}-\sqrt{3}}{2} \Rightarrow (C)$

 Note: Since $C$ and $D$ are not labeled in the diagram, refer to solution 1 for the location of points $C$ and $D$.

 



\textbf{Solution 6}


                    \begin{figure}[H]
                        \centering
                        \includegraphics[width=0.5\textwidth]{Solution/AMC_10_10A_2008_25_Solution_Figure_6.png}
                    \end{figure}
                    

 Notice that $\overarc{AE}$ is $\frac16$ the circumference of the circle. Therefore, $\overline{AE}$ is the side length of an inscribed hexagon with side length $4$. $\triangle AFE$ is a right triangle with $\overline{AF}=\frac12$. The length of $\overline{EF}=x+\frac{\sqrt{3}}{2}$. Using the \href{https://artofproblemsolving.com/wiki/index.php?title=Pythagorean\_Theorem}{Pythagorean Theorem}, we get

 
$4^2 = \left(\frac{1}{2}\right)^2 + \left(x+\frac{\sqrt{3}}{2}\right)^2$

 
Solving for $x$, we get $x = \frac{3\sqrt{7}-\sqrt{3}}{2}\ \boxed{\text{C}}$

 



\textbf{Solution 7 (Answer Choices & Estimation)}

The smallest distance between the intersection between two adjacent placemats and the circle is a bit less than $1$. Thus, the answer will be a bit more than $4-1=3.$ Going up, we guess that the next one up is the answer, $\frac{3\sqrt{7}-\sqrt{3}}{2},$ which is about $3.1$. The next one up after that is $2\sqrt{3},$ which is about $3.5$ and seems too big. So we circle $\boxed{\text{(C) }\frac{3\sqrt{7}-\sqrt{3}}{2}}$ and are happy we just cheesed the $AMC$ $10$ $2008$ problem $25.$

 


\end{enumerate}